% vim: ai sw=2

\documentclass[xcolor=dvipsnames,handout]{beamer}


\let\Oldemph\emph
\renewcommand{\emph}[1]{\textcolor{blue}{\Oldemph{#1}}}
\newcommand{\heading}[1]{\textbf{\textcolor{BurntOrange}{#1}}}
\newcommand{\header}[1]{\textbf{\texttt{\textcolor{WildStrawberry}{#1}}}}
\newcommand{\doc}[1]{\textbf{\texttt{\textcolor{OliveGreen}{#1}}}}
\newcommand{\pozor}{\emph{POZOR}}
\newcommand{\qtc}{\textsc{Qt Creator}}
\newcommand{\qtcu}{\textsc{Qt Creatoru}}

\mode<presentation>
{
\usetheme{Madrid}
}

\setbeamertemplate{navigation symbols}{}

\usepackage[czech]{babel}
\usepackage[utf8x]{inputenc}
\usepackage{helvet}
\usepackage{mathptmx}
\usepackage[T1]{fontenc}
\usepackage{graphicx}
\usepackage{hyperref}
\usepackage[final]{listings}
\usepackage{multirow}
\usepackage{tikz}
\usepackage[absolute,overlay]{textpos}
\usepackage{xspace}
\usetikzlibrary{arrows, fit, positioning}

\uselanguage{czech}
\deftranslation[to=czech]{Section}{Sekce}
\deftranslation[to=czech]{Subsection}{Podsekce}

\tikzset{tl_state/.style={rounded corners, thick, draw, minimum height=6mm,
	minimum width=12mm, bottom color=black!15, top color=black!5}}
\tikzset{tl_phase/.style={thick, draw, minimum size=8mm}}
\tikzset{tl_edge/.style={<-, >=stealth', semithick}}

\lstset{tabsize=2,numbers=none,showstringspaces=false,breaklines=true,
	breakatwhitespace=true,escapeinside={(*@}{@*)},
	basicstyle=\scriptsize,keywordstyle=\bfseries\color{OliveGreen},
	commentstyle=\color{blue}, columns=flexible, language=C,
	morekeywords={inline}}


\newcommand{\code}[2][]{\lstinline[columns=fullflexible, basicstyle=\ttfamily,
	keepspaces=true, #1]$#2$}
\newcommand{\cmd}[2][]{\code[language=bash, #1]{#2}}
\newcommand{\file}[2][]{\code[language={}, #1]{#2}}

\newcommand{\xes}{\texttt{x86}\xspace}

\newif\iffimuni
\fimunifalse

\newcommand{\mysection}[1]{\subsection*{#1}\frame{\subsectionpage}}

\title{Programování systému UNIX/Linux}

\author{Jiri~Slaby}

\institute[PC-DIR]{\includegraphics[trim=6cm 15cm 3.5cm 6cm,clip,width=2.5cm]{PCDIR_logo}}
